\section{Track I: Technical Session 4 - Quasi-Monte Carlo, Part I}\newpage

\begin{talk}
  {Halton Sequences, Scrambling and the Inverse Star-Discrepancy}% [1] talk title
  {Christian Wei\ss{}}% [2] speaker name
  {Ruhr West University of Applied Sciences}% [3] affiliations
  {christian.weiss@hs-ruhrwest.de}% [4] email
  {}% [5] coauthors
  
    
   
Halton sequences are classical examples of multi-dimensional low-discrepancy sequences. Braaten and Weller discovered that scrambling strongly reduces their empirical star-discre-pancy. A similar approach may be applied to certain multi-parameter subsequences of Halton sequences. Indeed, results from p-adic analysis guarantee that these subsequences still have the theoretical low-discrepancy property while scrambling has strong effects on the empirical star-discrepancy. By optimizing the parameters of these subsequences known empiric bounds for the inverse star-discrepancy can be improved.

\medskip


\begin{enumerate}
 \item[{[1]}] E. Braaten, and G. Weller. (1979). An improved low-discrepancy sequence for multidimensional quasi-Monte Carlo integration. Journal of Computational Physics, 33(2): 249--258.
 \item[{[2]}] C. Wei\ss{}. (2024). Scrambled Halton Subsequences and Inverse Star-Discrepancy, arXiv: 2411.10363. 
\end{enumerate}

\end{talk}

\begin{talk}
  {Transport Quasi-Monte Carlo}% [1] talk title
  {Sifan Liu}% [2] speaker name
  {Center for Computational Mathematics, Flatiron Institute}% [3] affiliations
  {sliu@flatironinstitute.org}% [4] email
  {}% [5] coauthors
  
    
   
Quasi-Monte Carlo (QMC) is a powerful method for evaluating high-dimensional integrals. However, its use is typically limited to distributions where direct sampling is straightforward, such as the uniform distribution on the unit hypercube or the Gaussian distribution. For general target distributions with potentially unnormalized densities, leveraging the low-discrepancy property of QMC to improve accuracy remains challenging. We propose training a transport map to push forward the uniform distribution on the unit hypercube to approximate the target distribution. Inspired by normalizing flows, the transport map is constructed as a composition of simple, invertible transformations. To ensure that RQMC achieves its superior error rate, the transport map must satisfy specific regularity conditions. We introduce a flexible parametrization for the transport map that not only meets these conditions but is also expressive enough to model complex distributions. Our theoretical analysis establishes that the proposed transport QMC estimator achieves faster convergence rates than standard Monte Carlo, under mild and easily verifiable growth conditions on the integrand. Numerical experiments confirm the theoretical results, demonstrating the effectiveness of the proposed method in Bayesian inference tasks.


\medskip

% If you would like to include references, please do so by creating a simple list numbered by [1], [2], [3], \ldots. See example below.
% Please do not use the \texttt{bibliography} environment or \texttt{bibtex} files.
% APA reference style is recommended.
% \begin{enumerate}
%  \item[{[1]}] Niederreiter, Harald (1992). {\it Random number generation and quasi-Monte Carlo methods}. Society for Industrial and Applied Mathematics (SIAM).
%  \item[{[2]}] Roberts, Gareth O, \& Rosenthal, Jeffrey S. (2002).  Optimal scaling for various Metropolis-Hastings algorithms, \textbf{16}(4), 351--367.
% \end{enumerate}

% Equations may be used if they are referenced. Please note that the equation numbers may be different (but will be cross-referenced correctly) in the final program book.
\end{talk}

\begin{talk}
  {Using Normalizing Flows for Efficient Quasi-Random Sampling for Copulas}% [1] talk title
  {Ambrose Emmett-Iwaniw}% [2] speaker name
  {University of Waterloo Department of Actuarial Science and Statistics}% [3] affiliations
  {arsemmettiwaniw@uwaterloo.ca}% [4] email
  {Christiane Lemieux}% [5] coauthors
  
    
   
In finance and risk management, copulas are used to model the dependence between stock prices and insurance losses to compute expectations of interest. Generally, Monte Carlo (MC) sampling is used to generate copula samples to approximate expectations. To reduce the variance of the approximation, we can use quasi-Monte Carlo (QMC) sampling to generate copula samples. This paper examines a new method to generate quasi-random samples from copulas requiring fewer training resources than previous methods such as the generative moment matching networks (GMMN) model [1]. Traditional methods that do not use generative models often rely on conditional distribution methods (CDM) to generate quasi-random samples from specific copulas [2]. CDM is limited to only a few parametric copulas (Gumbel has no efficient CDM to sample quasi-random samples) in low dimensions [2]. Here, we propose using a powerful and simple generative model called Normalizing Flows (NFs) to generate quasi-random samples for any copula, including cases where we only have data available. NFs are a type of explicit generative model that relies on transforming a simple density, such as a normal density, through efficient invertible transformations that rely on the change of variables formula into a density that models complex data that facilitates easy sampling and efficient inverting of samples from complex data to normal data and vice versa. The benefit of these NFs for copula modelling is that their training is efficient in terms of runtime, allowing for larger batch sizes compared to the GMMN model [1]. Also, it is sample-efficient; it only needs samples from the copula and not samples from the normal as the GMMN model [1] required. Once the NF model is trained, we can efficiently invert the model to take as input quasi-random samples to generate quasi-random copula samples. Through many different simulations and applications, we show our approach allows us to leverage the benefit of QMC in a variety of real-world settings involving dependent data.
\medskip

\begin{enumerate}
 \item[{[1]}] Hofert, M., Prasad, A., and Zhu, M. (2021). Quasi-random sampling for multivariate distributions via generative neural networks. Journal of Computational and Graphical Statistics,
30(3):647--670.
 \item[{[2]}] Cambou, M., Hofert, M., and Lemieux, C. (2017). Quasi-random numbers for copula models. Statistics and Computing, 27:1307--1329. 
\end{enumerate}
\end{talk}

\begin{talk}
  {Optimization of Kronecker Sequences}% [1] talk title
  {Claude D. Hall Jr.}% [2] speaker name
  {Illinois Institute of Technology}% [3] affiliations
  {cdhjr2000@gmail.com, chall14@iit.edu}% [4] email
  {Dr. Fred Hickernell, Dr. Nathan Kirk}% [5] coauthors
  
    

Our goal is to optimize Kronecker sequences [1,2] for Quasi-Monte Carlo integration. A shifted Kronecker sequence can be defined as $\{ x_i = i \boldsymbol{\alpha} + \boldsymbol{\Delta} \mod 1: i = 0, 1, 2, 3, ... \} \subset [0,1]^d $, where  $\boldsymbol{\alpha} \in [0,1)^d$ is our generating vector and $\boldsymbol{\Delta} \in [0,1)^d$ is the shift. There is no preferred sample size for Kronecker sequences, whereas the similarly defined lattice sequences do have preferred sample sizes. 

For discrepancies defined in terms of certain families of periodic reproducing kernels, when you replace any one of the elements $\alpha_j$ by $1 - \alpha_j$, the discrepancy stays unchanged. This allows one to shrink the search space for good generating vectors.  Performing a component by component search for $\boldsymbol{\alpha}$, we find that discrepancy for the Kronecker sequence converges nearly like $O(n^{-1})$, especially for smaller dimensions.  This talk presents the results of our component-by-component search algorithm and the performance of Kronecker sequences as measured by discrepancy and test cases.

\medskip

\begin{enumerate}
 \item[{[1]}] G.~Larcher, \emph{On the distribution of $s$-dimensional {K}ronecker-sequences}, Acta Arith. (1988), 335--347.
 \item[{[2]}] H.~Niederreiter, \emph{Quasi-Monte Carlo methods and pseudo-random numbers}, Bull. Amer. Math. Soc. 84 (1978), 957--1041.
\end{enumerate}
\end{talk}
